\documentclass[10pt,a4paper]{amsart}
\usepackage[margin=19mm]{geometry}
\usepackage{amsmath,amssymb,mathtools}
\usepackage[scr=rsfs,cal=euler]{mathalfa}
\usepackage{xfrac}
\usepackage[unicode,hidelinks,bookmarks=false]{hyperref}
\newcommand{\ie}{i.e.\ }
\newcommand{\cf}{cf.\ }
\newcommand{\resp}{respectively\ }
\newcommand{\Subsection}{\subsection}
\providecommand{\nobreakdash}{\nobreak\hbox{-}}
\setlength{\emergencystretch}{2em}
\multlinegap0pt
\usepackage{bm}
\usepackage{bbm}
\usepackage{dsfont}
\usepackage{xspace}
\usepackage{cancel}
\usepackage{mathtools}
\usepackage{amsthm}
\makeatletter
\@ifundefined{prop}{\newtheorem{prop}{Prop}}{}
\@ifundefined{theorem}{\newtheorem{theorem}[prop]{Theorem}}{}
\makeatother
\makeatletter
\expandafter\ifx\csname Rem\endcsname\relax
\newenvironment{Rem}{{\bf Remark.}}{}
\fi
\newcommand{\mymod}[1]{(\operatorname{mod} #1)}
\newcommand{\norm}{\mathrm{N}}
\makeatother
\title[Towards Keating-Snaith's conjecture for cubic Hecke $L$-func]{Towards Keating-Snaith's conjecture for cubic Hecke $L$-functions over the Eisenstein field}
\author{Hua Lin, Peng-Jie Wong}
\date{Source: arXiv:2511.08783v1 (CC BY 4.0)}
\begin{document}
\maketitle

\noindent\emph{Source and licence.} Towards Keating-Snaith's conjecture for cubic Hecke $L$-functions over the Eisenstein field. Version arXiv:2511.08783v1 (CC BY 4.0), distributed under the Creative Commons Attribution 4.0 International licence, as recorded in the arXiv licence metadata for this exact version and shown on the abstract page https://arxiv.org/abs/2511.08783v1. Full rights notice: licenses/arxiv-2511.08783-CC-BY-4.0.txt.\par\smallskip

\noindent\textit{Editorial scope.} Counted unit: the Theorem whose source label is \texttt{theorem S\_1} in the article (source0.tex lines 402--416, source label \texttt{theorem S\_1}), delivered together with the article's own complete original proof of it (source0.tex lines 419--485, one \texttt{proof} environment of 67 lines). No mathematics is written, repaired or re-derived: statement and proof are the article's own text line for line. Required context retained uncounted: none. Every definition, display and label of those lines is retained unchanged. Omitted from this package and not redistributed: 1--401, 417--418, 486--2034. Nothing inside a retained excerpt is omitted or altered: every retained body is delivered exactly as the article's lines are written, so no omission receipt of any kind accompanies this package. Citations: the retained excerpts carry no citation command at all, so no bibliography entry of the article is transcribed and no reference is redistributed. Reference adaptation: none was needed and none was made; every reference inside the delivered unit resolves inside it, so no unresolved reference or citation is produced. Printed slips kept literal: no printed word, symbol, hypothesis or number of the counted statement or of its proof was altered. Layout and class: the article's own class is \texttt{amsart} with packages including amssymb, amsmath, eucal, epstopdf, mathrsfs, enumerate, fullpage, setspace, xcolor, dsfont, todonotes, hyperref, fontenc, microtype; the wrapper is the lane's pinned \texttt{amsart} editorial wrapper at 10pt on A4, which restates the theorem-like environments the retained text needs and adds the article's own macro definitions that the retained text needs (\texttt{\textbackslash mymod}, \texttt{\textbackslash norm}), with extra wrapper packages (bm, bbm, dsfont, xspace, cancel, mathtools). The delivered numbering is the unit's own. Assets: no retained span contains a figure environment, a graphics-inclusion command, a PDF-page inclusion, a table or a TikZ picture; no image file is copied into this package, the package declares no asset dependency and no image file is redistributed with it. Licence: version arXiv:2511.08783v1 of this article is distributed under the Creative Commons Attribution 4.0 International licence, as recorded in the arXiv licence metadata for this exact version and shown on the abstract page https://arxiv.org/abs/2511.08783v1. A full rights notice is delivered at licenses/arxiv-2511.08783-CC-BY-4.0.txt. Characters: every retained excerpt is pure ASCII, so the delivered engine, XeLaTeX with its default Latin Modern text font, reports no missing character.
\begin{theorem}
\label{theorem S_1}
    Let $\ell\in\mathbb{Z}[\omega]$, such that $\norm(\ell)^{1/4}\le \sqrt{X}$, and let
    $$S=\sum_{\mathfrak{f}\in \mathcal{F}}\chi_{\mathfrak{f}}(\ell)\Phi\left(\frac{\norm(\mathfrak{f})}{X}\right),$$
    where $\Phi$ is a smooth function defined as above.   
    When $\ell$ is a cube, for any $A>0,$ we have
    $$S=\frac{X\widehat{\Phi}(0)}{\norm(9)}  \zeta^{-1}_\mathcal{K}(2)  \prod_{p\mid \ell}\left(1+\frac{1}{\norm(p)}\right)^{-1} +O\left(\frac{X}{A}\right)+O\left(\frac{A}{X^{1+\epsilon}}\right).$$
    When $\ell$ is not a cube, we have
    $$S
    \ll
    \frac{X}{A}+\frac{A}{X^{1+\epsilon}}.
    $$
Consequently, with the choice of $A=X^{1/2}/\norm(\ell)^{1/4}$, all the errors above are
$\ll X^{1/2}\norm(\ell)^{1/4}$.
\end{theorem}

\begin{proof} For brevity, let $|\cdot|$ denote the ideal norm in the following proof. Using 
\begin{equation}\label{detect-square}
\sum_{\substack{d\in \mathbb{Z}[\omega]\\d^2\mid \mathfrak{f}}}\mu_\mathcal{K}(d)=\begin{cases}
    1 &\text{ if } \mathfrak{f} \text{ is square-free,}\\
    0 &\text{ otherwise,}
\end{cases}
\end{equation}
to detect square-free $\mathfrak{f}$, and writing $\mathfrak{f}=d^2f$, we have
    \begin{align*}
    S
%    =
%    \sum_{\mathfrak{f}\in \mathcal{F}}\chi_{\mathfrak{f}}(\ell)\Phi\left(\frac{|\mathfrak{f}|}{X}\right)
    =
    \sum_{\substack{d\in \mathbb{Z}[\omega]}}\mu_\mathcal{K}(d)\chi_{d^2}(\ell)
    \sum_{f \equiv \overline{d^2}\mymod9}\chi_{f}(\ell)\Phi\left(\frac{|fd^2|}{X}\right).
\end{align*}

We trivially bound the sum for $|d|>A$ as
\begin{align*}
    \sum_{\substack{|d|>A}}\mu_\mathcal{K}(d)\chi_{d^2}(\ell)
    \sum_{f \equiv \overline{d^2}\mymod9}\chi_{f}(\ell)\Phi\left(\frac{|fd^2|}{X}\right)
    \ll
    \sum_{\substack{|d|>A}}\frac{X}{|d^2|}
    \ll
    \frac{X}{A}.
\end{align*}
For $|d|\le A$, after applying Poisson summation to the sum over $f$,
\begin{align}
\begin{split}
    &\sum_{\substack{|d|\le A}}\mu_\mathcal{K}(d)\chi_{d^2}(\ell)
    \sum_{f \equiv \overline{d^2}\mymod9}\chi_{f}(\ell)\Phi\left(\frac{|fd^2|}{X}\right)
    \\
    &=
    \frac{X}{|9|}\sum_{\substack{|d|\le A}}\frac{\mu_\mathcal{K}(d)\chi_{d^2}(\ell)}{|d^2|}
    \sum_{k\in \mathbb{Z}[\omega]} \widehat{\Phi}\bigg(\sqrt{\frac{X|k|}{|9{\ell }d^2|}}\bigg)
    \sum_{r\mymod {\ell }} \frac{\chi_{\ell }(r)}{|\ell|}e\bigg(\frac{-( r9\overline{9}+\overline{d^2}\ell\overline{\ell})k}{{9\ell }\sqrt{-3}}\bigg).
\end{split}
\end{align}
When $k=0$, the above contribution is nonzero if and only if $\ell$ is a cube, and in that case
$$S=\frac{X\widehat{\Phi}(0)}{|9|}  \zeta^{-1}_\mathcal{K}(2)  \prod_{p\mid \ell}\left(1+\frac{1}{|p|}\right)^{-1} +O\left(\frac{X}{A}\right).$$
When $k\ne 0,$ we have
\begin{align*}
    &\frac{X}{|9|}\sum_{\substack{|d|\le A\\3\nmid d}}\frac{\mu_\mathcal{K}(d)\chi_{d^2}(\ell)}{|d^2|}
    \sum_{\substack{k\in \mathbb{Z}[\omega]\\k\ne 0}} \widehat{\Phi}\bigg(\sqrt{\frac{X|k|}{|9{\ell }d^2|}}\bigg)
    e\bigg(\frac{- \overline{d^2\ell}k}{{9 }\sqrt{-3}}\bigg)
    \frac{1}{|\ell|}\sum_{r\mymod {\ell }} \chi_{\ell }(r)e\left(\frac{- r\overline{9}k}{{\ell }\sqrt{-3}}\right)\\
    &\ll
    \sum_{\substack{|d|\le A}}\frac{X}{|9d^2|} \frac{|9\ell d^2|}{X} \left(\frac{|9\ell d^2|}{X}\right)^{K-1}
    \sum_{\substack{k\in \mathbb{Z}[\omega]\\k\ne 0}}
    \left(\frac{1}{|k|}\right)^K
    \ll
    \frac{|9|^{K-1}|\ell|}{X^{\epsilon K}}\sum_{\substack{|d|\le A}} 
    \sum_{\substack{k\in \mathbb{Z}[\omega]\\k\ne 0}}
    \left(\frac{1}{|k|}\right)^K\ll
    \frac{|\ell|A}{X},
\end{align*}
where since
$$\frac{X}{|\ell d^2|} \ge X^\epsilon \text{ for some } \epsilon<1, \text{ and } \  \widehat{\Phi}(x)\ll_K x^{-K} \text{ for any } K>0,$$
we choose $K>1$ such that $\epsilon K=1.$
Thus, when $\ell$ is a cube, we have 
$$S=\frac{X\widehat{\Phi}(0)}{|9|}  \zeta^{-1}_\mathcal{K}(2)  \prod_{p\mid \ell}\left(1+\frac{1}{|p|}\right)^{-1} +O\left(\frac{X}{A}\right)+O\left(\frac{|\ell|A}{X}\right),$$
and when $\ell$ is not a cube, 
$$S
\ll
\frac{X}{A}+\frac{|\ell|A}{X},$$
which completes the proof.
\end{proof}

\end{document}
